\textbf{Efficiency: avoiding repeated work}

A cache stores results so a later request can reuse them. Memoization is
caching a function's return value for each input. On a \emph{cache miss},
compute and store the result; on a \emph{cache hit}, return the stored result.

\begin{lstlisting}
def memo(f):
    cache = {}
    def remembered(x):
        if x not in cache:
            cache[x] = f(x)
        return cache[x]
    return remembered
\end{lstlisting}

Each call to \lstinline{memo} creates a new dictionary. The returned function
keeps that dictionary across calls. Putting \lstinline|cache = {}| inside
\lstinline{remembered} would discard earlier results on every call.
Use membership testing: a stored result of 0 or \lstinline{False} is still
a cached result. This version requires hashable inputs, such as integers.

Caching works when the same input should produce the same result. A cache
hit skips the original function body, including its printing or other side
effects. If a result depends on changing external state, reusing it can be
incorrect. Caching also uses memory to store results.

\textbf{Fibonacci and repeated calls}

Tree-recursive Fibonacci repeats calls on smaller inputs.
For example, computing \lstinline{fib(5)} needs the results for inputs 4
and 3. Computing the result for input 4 needs input 3 again.
With memoization, every recursive call checks a shared dictionary keyed by
the argument. Each \lstinline{fib(k)} is computed only once; later requests
are just lookups. Saving only the final answer to the outermost call would
not prevent repeated work within that computation.

\textbf{List operations differ in cost}
\begin{itemize}
    \item \textbf{Fast list-building choices:} \lstinline{s.append(x)} adds
    to the end without shifting existing elements. List comprehensions also
    build results efficiently, evaluating their expression once per element.
    Appending occasionally needs more storage, but does not copy the whole
    list every time.
    \item \textbf{Fast replacement:} \lstinline{s[k] = k * k} replaces
    the element at an existing index without shifting other elements or
    copying the list. It does not grow the list: index \lstinline{k} must
    already be valid.
    \item \textbf{Slow when repeated:} \lstinline{s.insert(0, x)} shifts
    existing elements to make room at the front. \lstinline{s = s + [x]}
    creates a new list, copying the growing old list every time.
\end{itemize}
Explain what work is repeated or avoided, rather than giving a runtime formula.
